\section*{Capítulo \ref{ch:b3:structural-analysis} --- Análisis estructural}

\begin{solution}{exo:b3:structural-analysis:1}
El \omterm{def:b3:structural-analysis:stereo}{polo} tiene una dirección de $40 + 180 = \ang{220}$ y una inmersión de $90 - 30 = \ang{60}$. El
plano buza $90 - 20 = \ang{70}$ hacia $315 - 180 = \ang{135}$: $135/70$.
\end{solution}

\begin{solution}{exo:b3:structural-analysis:2}
$\theta = \ang{60}$. \omterm{def:b3:structural-analysis:stereo}{Red equiangular}: $10\tan\ang{30} = \qty{5.77}{cm}$;
en la red equivalente: $10\sqrt2\sin\ang{30} = \qty{7.07}{cm}$.
\end{solution}

\begin{solution}{exo:b3:structural-analysis:3}
$\vec n_1 = (\sin\ang{30}, 0, \cos\ang{30})$, $\vec n_2 = (-\sin\ang{30}, 0,
\cos\ang{30})$, $\vec n_1\cdot\vec n_2 = -0.25 + 0.75 = 0.5$: los \omterm{def:b3:structural-analysis:stereo}{polos} están
separados \ang{60}, el ángulo entre flancos es \ang{120}, un pliegue abierto.
\end{solution}

\begin{solution}{exo:b3:structural-analysis:4}
$\vec d = (\cos\ang{60}, 0, -\sin\ang{60}) = (0.5, 0, -0.866)$ y $\vec u =
0.707\,\vec s + 0.707\,\vec d = (0.354, 0.707, -0.612)$: dirección
$\operatorname{atan2}(0.354, 0.707) = \ang{26.6}$, inmersión $\arcsin0.612 =
\ang{37.8}$. El cabeceo se sitúa entre \ang{-90} (normal) y 0 (sinestral): una
falla normal oblicua con componente sinestral.
\end{solution}

\begin{solution}{exo:b3:structural-analysis:5}
\begin{gather*}
  \vec n_1 = (0.5567, 0.3214, 0.7660),\qquad \vec n_2 = (-0.4698, -0.1710, 0.8660),\\
  \vec n_1\times\vec n_2 = (0.4093, -0.8420, 0.0558),
\end{gather*}
\sloppy de longitud $0.9379$. Su opuesto, $-\vec n_1\times\vec n_2$, apunta hacia abajo:
dirección \ang{334.1}, inmersión $\arcsin(0.0558/0.9379) = \ang{3.4}$. $\vec n_1\cdot\vec
n_2 = 0.3469$, \ang{69.7} entre las normales, ángulo entre flancos \ang{110.3}:
un pliegue abierto.
\end{solution}

\begin{solution}{exo:b3:structural-analysis:6}
Una inmersión superior a \ang{60} significa $\theta < \ang{30}$. En la red equivalente: $r^2 =
2\sin^2\ang{15} = 0.134$, \qty{13.4}{\percent} de la red; en la \omterm{def:b3:structural-analysis:stereo}{red equiangular}:
$\tan^2\ang{15} = 0.072$, \qty{7.2}{\percent}. En la esfera el casquete contiene
$1 - \cos\ang{30} = 0.134$ del hemisferio: solo la red equivalente lo muestra
con su proporción real, por eso los \omterm{def:b3:structural-analysis:stereo}{polos} se cuentan sobre ella.
\end{solution}

\begin{solution}{exo:b3:structural-analysis:7}
\begin{gather*}
  T_{xx} = \frac{2\sin^2\ang{20}}3 = 0.078,\qquad
  T_{zz} = \frac{2\cos^2\ang{20} + 1}3 = 0.922,\\
  T_{xz} = -\frac{\sin\ang{20}\cos\ang{20} - \sin\ang{20}\cos\ang{20}}3 = 0,
\end{gather*}
y todo término con $y$ es 0. Los autovalores son $0.922$ (vertical),
$0.078$ (este) y $0$ (norte): el \omterm{def:b3:structural-analysis:pi}{eje beta} es horizontal, norte--sur.
Los tres \omterm{def:b3:structural-analysis:stereo}{polos} están en un mismo plano, así que ninguno tiene componente según
su normal, y $\vec u^{\mathsf T}T\vec u = 0$ para $\vec u$ según ella.
\end{solution}

\begin{solution}{exo:b3:structural-analysis:8}
$a = 2$, $b = 1.5$, $k = 1/0.5 = 2$: constricción. $Z = 1/(1.8\times1.1) =
0.505$, $a = 1.636$, $b = 2.178$, $k = 0.636/1.178 = 0.54$: achatamiento.
\end{solution}

\begin{solution}{exo:b3:structural-analysis:9}
$S_1/S_2 = 2.25$ y $S_1S_2 = 1$: $S_1 = 1.5$, $S_2 = 0.667$; el eje largo
se alargó un \qty{50}{\percent}. El poro tiene los semiejes $1.5\times0.4 =
\qty{0.60}{mm}$ y $0.667\times0.4 = \qty{0.27}{mm}$.
\end{solution}

\begin{solution}{exo:b3:structural-analysis:10}
\sloppy \omterm{def:b3:structural-analysis:stereo}{Red equiangular}: $1/(2\cos^2\ang{30}) = 0.667$ a lo largo del radio y
$\tan\ang{30}/\sin\ang{60} = 0.577/0.866 = 0.667$ alrededor: iguales, conformes.
Red equivalente: $\mathrm dr/(R\,\mathrm d\theta) = (\sqrt2/2)\cos\ang{30} = 0.612$
y $\sqrt2\sin\ang{30}/\sin\ang{60} = 0.707/0.866 = 0.816$; el producto es
$0.500$, el factor de área del teorema. Un círculo menor se convierte en un óvalo
de ejes en la razón $0.816/0.612 = 1.33$, largo alrededor del centro.
\end{solution}

\begin{solution}{exo:b3:structural-analysis:11}
La falla es vertical, así que su normal es horizontal y perpendicular al
$\sigma_2$ vertical; por la parte (b) del teorema el deslizamiento es perpendicular a
$\sigma_2$, horizontal, para cualquier $\Phi$. Con $\sigma_1 = 1$ según el norte,
$\sigma_3 = 0$ según el este: $S\vec n = (0, -0.342, 0)$, $\vec n\cdot S\vec n =
0.117$, $\vec\tau = (0, -0.342, 0) - 0.117\,(0.940, -0.342, 0) = (-0.110,
-0.302, 0)$ y $-\vec\tau = (0.110, 0.302, 0)$, hacia \ang{020}. El bloque
en el lado de $\vec n$, el lado este, se desplaza hacia el norte-noreste a lo largo
del rumbo: visto desde el lado oeste, se mueve hacia la izquierda, una falla sinestral (cabeceo
0). $\sigma_2$ no interviene en el cálculo.
\end{solution}

\begin{solution}{exo:b3:structural-analysis:12}
$S = (7.5 - 6.0)/1.0 = \qty{1.5}{km}$; $h = 6.0/1.5 = \qty{4.0}{km}$ por debajo del
horizonte inferior. La sección tenía $20 + 1.5 = \qty{21.5}{km}$ de longitud, y cambió
en $-1.5/21.5 = \qty{-7.0}{\percent}$.
\end{solution}

\begin{solution}{pb:b3:structural-analysis:1}
\textbf{1.} El \omterm{def:b3:structural-analysis:stereo}{polo} tiene una dirección de \ang{255} y una inmersión de \ang{77}; $\theta =
\ang{13}$ y $r = 10\sqrt2\sin\ang{6.5} = \qty{1.60}{cm}$ desde el centro,
hacia \ang{255}.

\textbf{2.} $\vec n = (\sin\ang{13}\sin\ang{75}, \sin\ang{13}\cos\ang{75},
\cos\ang{13}) = (0.2173, 0.0582, 0.9744)$. Seis vectores unitarios casi paralelos
suman casi 6: $|\vec S_E| = 5.9974$.

\textbf{3.} Flanco este: $\bar\alpha = \operatorname{atan2}(1.2583, 0.2017) =
\ang{80.9}$ y $\bar\delta = \arccos(5.8604/5.9974) = \ang{12.3}$; $\bar R
= 5.9974/6 = 0.99956$ y la dispersión es $\sqrt{2\times0.00044} =
\qty{0.0297}{rad}$, o \ang{1.7}. Flanco oeste: $\bar\alpha =
\operatorname{atan2}(-1.2445, 0.1886) = \ang{278.6}$ y $\bar\delta =
\arccos(5.8636/5.9972) = \ang{12.1}$; $\bar R = 0.99953$, dispersión \ang{1.8}.

\textbf{4.}
\[
  \vec S_E\times\vec S_W = (0.0774, -14.6714, 0.4883),
\]
de longitud $14.680$. Su opuesto apunta hacia abajo: la dirección es $\operatorname{atan2}(-0.0774, 14.6714) = \ang{359.7}$ y la
inmersión $\arcsin(0.4883/14.680) = \ang{1.9}$.

\textbf{5.} El coseno del ángulo entre las normales medias es
$32.835/35.968 = 0.9129$, el producto escalar de las sumas dividido por sus
longitudes: \ang{24.1} entre las normales, un ángulo entre flancos de \ang{155.9}:
un pliegue abierto, cercano a suave.

\textbf{6.}
\begin{gather*}
  \vec m_E = (0.2098, 0.0336, 0.9772),\qquad \vec m_W = (-0.2075, 0.0314, 0.9777),\\
  \vec m_E - \vec m_W = (0.4173, 0.0022, -0.0005):
\end{gather*}
una normal horizontal hacia \ang{089.7}. El plano axial es vertical, con rumbo
\ang{359.7}, prácticamente norte--sur: el pliegue es recto.

\textbf{7.} El cinturón es el plano cuyo \omterm{def:b3:structural-analysis:stereo}{polo} es el \omterm{def:b3:structural-analysis:pi}{eje beta}: dirección de buzamiento
\ang{179.7}, buzamiento $90 - 1.9 = \ang{88.1}$, casi el plano vertical este--oeste.
Cada flanco contiene el eje de pliegue, así que cada normal es perpendicular a
él y yace en el plano perpendicular a él.

\textbf{8.} $0.04400 + 0.00144 + 0.95456 = 1.00000$. $\vec b^{\mathsf T}T\vec
b = 0.00038$, el menor autovalor: $\vec b$ es, a menos del redondeo, el
autovector $\vec e_3$, y el cinturón que pasa por los dos \omterm{def:b3:structural-analysis:stereo}{polos} medios es el
cinturón de mejor ajuste de los doce.

\textbf{9.} $\arcsin\sqrt{0.00038} = \arcsin0.0195 = \ang{1.1}$.

\textbf{10.} A lo largo del cinturón, los \omterm{def:b3:structural-analysis:stereo}{polos} se sitúan a unos \ang{12.2} de la vertical,
a uno y otro lado; $\sin^2\ang{12.2} = 0.045$, el cuadrado medio de sus
componentes este--oeste, que es $\lambda_2$.

\textbf{11.} Sí: los \omterm{def:b3:structural-analysis:stereo}{polos} se apartan del cinturón \ang{1.1} de media, no
más que la dispersión de \ang{1.7} de las lecturas de un flanco respecto de su media.

\textbf{12.} $\ang{1.1}/\sqrt{12} = \ang{0.3}$: la inmersión es \ang{1.9} $\pm$
\ang{0.3}.

\textbf{13.} Un cabeceo de \ang{-90}: una falla normal pura, con deslizamiento de buzamiento. Anderson:
$\sigma_1$ vertical, $\sigma_3$ horizontal según la dirección de buzamiento,
\ang{120}--\ang{300}, $\sigma_2$ horizontal según el rumbo,
\ang{030}--\ang{210}.

\textbf{14.} La normal de $120/60$ yace en el plano vertical de dirección
\ang{120}, que contiene $\vec e_1$ y $\vec e_3$: $n_2 = 0$. Por el apartado (b), la
deslizamiento es perpendicular a $\sigma_2$, el rumbo, para todo $\Phi$: es
deslizamiento de buzamiento. La falla no dice nada de $\Phi$; en la fórmula se anulan tanto $n_2$ como
$\vec e_2\cdot\vec w$, y por tanto el numerador y el denominador, se anulan.

\textbf{15.} $\vec n = (\sin\ang{55}\sin\ang{150}, \sin\ang{55}\cos\ang{150},
\cos\ang{55}) = (0.4096, -0.7094, 0.5736)$; $n_1 = 0.574$; $n_2 = 0.5\times
0.4096 + 0.866\times(-0.7094) = -0.410$.

\textbf{16.} $\vec e_1\cdot\vec w = 0.198$ y $\vec e_2\cdot\vec w = 0.5\times
0.771 + 0.866\times0.605 = 0.909$, así que $\Phi = -(0.574\times0.198)/(-0.410\times
0.909) = 0.1137/0.3727 = 0.305$.

\textbf{17.} $(0.305 + 0.294 + 0.301 + 0.306)/4 = 0.30$. Los cabeceos se
redondearon al grado; un grado de cabeceo mueve $\Phi$ alrededor de 0.02 en estas
fallas, así que el medio grado de redondeo explica la dispersión.

\textbf{18.} $\sigma_2 - \sigma_3 = 0.3(\sigma_1 - \sigma_3)$: $\sigma_2$ está
más cerca de $\sigma_3$ que de $\sigma_1$. Los dos esfuerzos horizontales están muy
por debajo del vertical; la región se distendió sobre todo hacia \ang{120}, un
poco según \ang{030}.

\textbf{19.} Con $\vec n = (0.750, -0.433, 0.500)$ y el deslizamiento $\vec u =
(0.433, -0.250, -0.866)$, $(\vec n - \vec u)/\sqrt2$ tiene una inmersión de \ang{75} hacia
\ang{300}: P; $(\vec n + \vec u)/\sqrt2$ tiene una inmersión de \ang{15} hacia \ang{120}:
T. Ambos están en el mismo plano vertical que $\sigma_1$ y $\sigma_3$, pero
\ang{15} de ellos, porque la falla forma \ang{30}, y no \ang{45}, con
$\sigma_1$.

\textbf{20.} $\tan a = \tan\ang{12.3}\times\cos(\ang{90} - \ang{80.9}) =
0.2180\times0.9874 = 0.2153$: $a = \ang{12.1}$.

\textbf{21.} $L_0 = 8000/\cos\ang{12} = \qty{8178.7}{m}$, $S =
\qty{178.7}{m}$, un cambio relativo de $\cos\ang{12} - 1 =
\qty{-2.2}{\percent}$.

\textbf{22.} $4000\tan\ang{12} = \qty{850.2}{m}$; $A = \tfrac12\times8000\times
850.2 = \qty{3.40e6}{m^2}$.

\textbf{23.} $h = \num{3.40e6}/178.7 = \qty{19.0}{km}$. Un despegue en la
base de la cobertera daría un $h$ de unos pocos cientos de metros y un área excedente
unas cincuenta veces menor: los pliegues no están despegados allí. O bien el
zócalo se plegó junto con la cobertera, o los sinclinales se hundieron por debajo del
nivel original, lo que hace $A$ menor; la sección por sí sola no puede decirlo.

\textbf{24.} Salto horizontal $200/\tan\ang{60} = \qty{115.5}{m}$, desplazamiento $200/\sin\ang{60}
= \qty{230.9}{m}$.

\textbf{25.} $\tan a = \tan\ang{60}\cos\ang{30} = 1.500$: $a = \ang{56.3}$, y
el salto horizontal aparente es $200/1.5 = \qty{133.3}{m}$. El desplazamiento de la falla es
según el buzamiento, hacia \ang{120}: la sección según \ang{120} muestra las
longitudes verdaderas y debe ser la que se restaure.

\textbf{26.} La falla corta los pliegues, así que es más joven y se restaura primero,
desplazando el bloque superior hacia arriba por la falla \qty{230.9}{m}; los pliegues se
desdoblan después. Las \omterm{def:b3:structural-analysis:restoration}{líneas de anclaje} van sobre las superficies axiales de los sinclinales, en el
bloque inferior, donde los pliegues simétricos no tenían deslizamiento entre las capas.

\textbf{27.} El eje de pliegue de Varden tiene una inmersión de
\textbf{\ang[text-series-to-math=true]{1.9}} hacia el norte (dirección
\ang{359.7}): los \ang{2} impuestos a la región, recuperados hasta el
\ang{0.3} que permiten las doce lecturas.
\end{solution}
